% Options for packages loaded elsewhere
\PassOptionsToPackage{unicode}{hyperref}
\PassOptionsToPackage{hyphens}{url}
%
\documentclass[
  10pt,
]{article}
\usepackage{amsmath,amssymb}
\usepackage{iftex}
\ifPDFTeX
  \usepackage[T1]{fontenc}
  \usepackage[utf8]{inputenc}
  \usepackage{textcomp} % provide euro and other symbols
\else % if luatex or xetex
  \usepackage{unicode-math} % this also loads fontspec
  \defaultfontfeatures{Scale=MatchLowercase}
  \defaultfontfeatures[\rmfamily]{Ligatures=TeX,Scale=1}
\fi
\usepackage{lmodern}
\ifPDFTeX\else
  % xetex/luatex font selection
\fi
% Use upquote if available, for straight quotes in verbatim environments
\IfFileExists{upquote.sty}{\usepackage{upquote}}{}
\IfFileExists{microtype.sty}{% use microtype if available
  \usepackage[]{microtype}
  \UseMicrotypeSet[protrusion]{basicmath} % disable protrusion for tt fonts
}{}
\makeatletter
\@ifundefined{KOMAClassName}{% if non-KOMA class
  \IfFileExists{parskip.sty}{%
    \usepackage{parskip}
  }{% else
    \setlength{\parindent}{0pt}
    \setlength{\parskip}{6pt plus 2pt minus 1pt}}
}{% if KOMA class
  \KOMAoptions{parskip=half}}
\makeatother
\usepackage{xcolor}
\usepackage[margin=1in]{geometry}
\usepackage{color}
\usepackage{fancyvrb}
\newcommand{\VerbBar}{|}
\newcommand{\VERB}{\Verb[commandchars=\\\{\}]}
\DefineVerbatimEnvironment{Highlighting}{Verbatim}{commandchars=\\\{\}}
% Add ',fontsize=\small' for more characters per line
\newenvironment{Shaded}{}{}
\newcommand{\AlertTok}[1]{\textcolor[rgb]{1.00,0.00,0.00}{\textbf{#1}}}
\newcommand{\AnnotationTok}[1]{\textcolor[rgb]{0.38,0.63,0.69}{\textbf{\textit{#1}}}}
\newcommand{\AttributeTok}[1]{\textcolor[rgb]{0.49,0.56,0.16}{#1}}
\newcommand{\BaseNTok}[1]{\textcolor[rgb]{0.25,0.63,0.44}{#1}}
\newcommand{\BuiltInTok}[1]{\textcolor[rgb]{0.00,0.50,0.00}{#1}}
\newcommand{\CharTok}[1]{\textcolor[rgb]{0.25,0.44,0.63}{#1}}
\newcommand{\CommentTok}[1]{\textcolor[rgb]{0.38,0.63,0.69}{\textit{#1}}}
\newcommand{\CommentVarTok}[1]{\textcolor[rgb]{0.38,0.63,0.69}{\textbf{\textit{#1}}}}
\newcommand{\ConstantTok}[1]{\textcolor[rgb]{0.53,0.00,0.00}{#1}}
\newcommand{\ControlFlowTok}[1]{\textcolor[rgb]{0.00,0.44,0.13}{\textbf{#1}}}
\newcommand{\DataTypeTok}[1]{\textcolor[rgb]{0.56,0.13,0.00}{#1}}
\newcommand{\DecValTok}[1]{\textcolor[rgb]{0.25,0.63,0.44}{#1}}
\newcommand{\DocumentationTok}[1]{\textcolor[rgb]{0.73,0.13,0.13}{\textit{#1}}}
\newcommand{\ErrorTok}[1]{\textcolor[rgb]{1.00,0.00,0.00}{\textbf{#1}}}
\newcommand{\ExtensionTok}[1]{#1}
\newcommand{\FloatTok}[1]{\textcolor[rgb]{0.25,0.63,0.44}{#1}}
\newcommand{\FunctionTok}[1]{\textcolor[rgb]{0.02,0.16,0.49}{#1}}
\newcommand{\ImportTok}[1]{\textcolor[rgb]{0.00,0.50,0.00}{\textbf{#1}}}
\newcommand{\InformationTok}[1]{\textcolor[rgb]{0.38,0.63,0.69}{\textbf{\textit{#1}}}}
\newcommand{\KeywordTok}[1]{\textcolor[rgb]{0.00,0.44,0.13}{\textbf{#1}}}
\newcommand{\NormalTok}[1]{#1}
\newcommand{\OperatorTok}[1]{\textcolor[rgb]{0.40,0.40,0.40}{#1}}
\newcommand{\OtherTok}[1]{\textcolor[rgb]{0.00,0.44,0.13}{#1}}
\newcommand{\PreprocessorTok}[1]{\textcolor[rgb]{0.74,0.48,0.00}{#1}}
\newcommand{\RegionMarkerTok}[1]{#1}
\newcommand{\SpecialCharTok}[1]{\textcolor[rgb]{0.25,0.44,0.63}{#1}}
\newcommand{\SpecialStringTok}[1]{\textcolor[rgb]{0.73,0.40,0.53}{#1}}
\newcommand{\StringTok}[1]{\textcolor[rgb]{0.25,0.44,0.63}{#1}}
\newcommand{\VariableTok}[1]{\textcolor[rgb]{0.10,0.09,0.49}{#1}}
\newcommand{\VerbatimStringTok}[1]{\textcolor[rgb]{0.25,0.44,0.63}{#1}}
\newcommand{\WarningTok}[1]{\textcolor[rgb]{0.38,0.63,0.69}{\textbf{\textit{#1}}}}
\usepackage{longtable,booktabs,array}
\usepackage{calc} % for calculating minipage widths
% Correct order of tables after \paragraph or \subparagraph
\usepackage{etoolbox}
\makeatletter
\patchcmd\longtable{\par}{\if@noskipsec\mbox{}\fi\par}{}{}
\makeatother
% Allow footnotes in longtable head/foot
\IfFileExists{footnotehyper.sty}{\usepackage{footnotehyper}}{\usepackage{footnote}}
\makesavenoteenv{longtable}
\setlength{\emergencystretch}{3em} % prevent overfull lines
\providecommand{\tightlist}{%
  \setlength{\itemsep}{0pt}\setlength{\parskip}{0pt}}
\setcounter{secnumdepth}{-\maxdimen} % remove section numbering
% definitions for citeproc citations
\NewDocumentCommand\citeproctext{}{}
\NewDocumentCommand\citeproc{mm}{%
  \begingroup\def\citeproctext{#2}\cite{#1}\endgroup}
\makeatletter
 % allow citations to break across lines
 \let\@cite@ofmt\@firstofone
 % avoid brackets around text for \cite:
 \def\@biblabel#1{}
 \def\@cite#1#2{{#1\if@tempswa , #2\fi}}
\makeatother
\newlength{\cslhangindent}
\setlength{\cslhangindent}{1.5em}
\newlength{\csllabelwidth}
\setlength{\csllabelwidth}{3em}
\newenvironment{CSLReferences}[2] % #1 hanging-indent, #2 entry-spacing
 {\begin{list}{}{%
  \setlength{\itemindent}{0pt}
  \setlength{\leftmargin}{0pt}
  \setlength{\parsep}{0pt}
  % turn on hanging indent if param 1 is 1
  \ifodd #1
   \setlength{\leftmargin}{\cslhangindent}
   \setlength{\itemindent}{-1\cslhangindent}
  \fi
  % set entry spacing
  \setlength{\itemsep}{#2\baselineskip}}}
 {\end{list}}
\usepackage{calc}
\newcommand{\CSLBlock}[1]{\hfill\break\parbox[t]{\linewidth}{\strut\ignorespaces#1\strut}}
\newcommand{\CSLLeftMargin}[1]{\parbox[t]{\csllabelwidth}{\strut#1\strut}}
\newcommand{\CSLRightInline}[1]{\parbox[t]{\linewidth - \csllabelwidth}{\strut#1\strut}}
\newcommand{\CSLIndent}[1]{\hspace{\cslhangindent}#1}
\ifLuaTeX
  \usepackage{selnolig}  % disable illegal ligatures
\fi
\usepackage{bookmark}
\IfFileExists{xurl.sty}{\usepackage{xurl}}{} % add URL line breaks if available
\urlstyle{same}
\hypersetup{
  pdftitle={Compiled Proof Guidance for Repository-Scale Verified Software Synthesis},
  pdfauthor={Mohit Mahajan (ByteVerity / Synthesa)},
  hidelinks,
  pdfcreator={LaTeX via pandoc}}

\title{Compiled Proof Guidance for Repository-Scale Verified Software
Synthesis}
\usepackage{etoolbox}
\makeatletter
\providecommand{\subtitle}[1]{% add subtitle to \maketitle
  \apptocmd{\@title}{\par {\large #1 \par}}{}{}
}
\makeatother
\subtitle{A Deterministic-Authority Study on Vero}
\author{Mohit Mahajan (ByteVerity / Synthesa)}
\date{September 2026}

\begin{document}
\maketitle
\begin{abstract}
Repository-scale verified software synthesis is usually approached as an
end-to-end agent task: a model edits implementations and proofs,
observes compiler feedback, and tries again. We study a different
architecture in which the model, when used, has no proof authority.
Synthesa resolves Vero specifications through Lean's elaborated
environment, compiles them into semantic obligation graphs, classifies
parent-contract adequacy, applies a bounded deterministic proof floor,
and promotes only kernel-checked helpers into a content-addressed
library. Exact residuals can request a narrowly typed proposal; the
response is pinned before Lean either admits or refutes it. Candidate
slot changes then pass through a signed, terminal-once build plan and
Vero's independent clean re-render and grader. In an iterative,
benchmark-seen campaign at Vero commit 0a7325d and Lean 4.29.1, the
resulting mixed-mode portfolio improves from 2/43 complete repositories
and 1030/2705 accepted specifications to 40/43 and 2606/2705, with no
Synthesa-only result counted as a Vero pass. The selected artifacts
comprise 37 proof-mode and six codeproof-mode evaluations, so 40/43 is
not presented as a uniform-mode leaderboard result. Thirty-three metered
external proposer calls cost USD 0.7104 in total; frozen artifact replay
makes zero model calls, while human, Codex-orchestration, and
local-compute costs were not fully metered. Three residual repositories
reveal an equally important result: fail-closed elaboration and
parent-adequacy checks distinguish proof-search exhaustion from
malformed specifications and missing semantic contracts. The evidence
supports compiled repository organization as a promising complement to
generative proof search, while motivating a future blind, uniform-mode
evaluation.
\end{abstract}

{
\setcounter{tocdepth}{3}
\tableofcontents
}
\section{1. Introduction}\label{introduction}

Machine-checked verification changes the authority structure of software
generation. A model can suggest an implementation or proof, but validity
is a property of a formal artifact checked by a small trusted kernel.
This makes verified synthesis a natural setting for separating creative
proposal from acceptance. The separation is especially important at
repository scale, where proof obligations share definitions, invariants,
and helper chains across modules.

Vero makes this repository-level problem concrete. It contains 43 Lean 4
projects, 743 APIs, and 2,705 formal specifications translated from real
software repositories. It supports proof-only evaluation against
supplied reference implementations and code-and-proof evaluation in
which implementations are also synthesized (\citeproc{ref-vero2026}{Ye,
Lou, et al. 2026}). Its strongest reported configuration completes 27
repositories in codeproof and 25 in proof mode, despite accepting more
than 85\% of individual specifications. Vero's analysis attributes this
gap to repository organization: completed repositories place a median of
more than 70\% of proof lines in helper lemmas, and almost every
successful repository shares helpers across specifications
(\citeproc{ref-vero2026}{Ye, Lou, et al. 2026}).

That observation suggests a systems question. Must a model discover the
entire repository proof architecture by free-form trial and error, or
can much of the organization be compiled? Synthesa explores the latter.
It treats repository verification as a collection of bounded decisions
connected by checked seams: which declaration a specification denotes,
whether its type elaborates, which implementation symbols and parent
laws it depends on, which goals share a semantic shape, which verified
abstraction covers them, and whether an exact candidate may be admitted.
The potentially unbounded creative step is narrowed to proposing a
missing abstraction only after deterministic decomposition.

We make four contributions:

\begin{enumerate}
\def\labelenumi{\arabic{enumi}.}
\tightlist
\item
  We describe an elaboration-first representation of repository proof
  obligations and a structured residual taxonomy that avoids source-text
  guessing.
\item
  We present a compiled proof-guidance architecture combining
  deterministic tactics, parent-contract adequacy, semantic clustering,
  precompiled formal oracles, verified helper reuse, and an optional
  bounded proposer.
\item
  We define an authority chain in which proposal bytes are pinned, Lean
  alone admits proofs, signed plans constrain materialization, and
  Vero's untouched grader remains the benchmark verdict.
\item
  We report a fully disclosed, iterative Vero campaign that produces a
  40/43 mixed-mode qualified portfolio and 2606/2705 accepted
  specifications, together with its costs, artifacts, and limitations.
\end{enumerate}

This is not a claim of a blind 40/43 single-mode agent run. It is a
systems study of whether repository-scale organization can be moved from
an agent into a deterministic proof-guided product.

\section{2. Problem formulation and trust
model}\label{problem-formulation-and-trust-model}

\subsection{2.1 Repository-level
obligations}\label{repository-level-obligations}

Let a Vero repository provide an implementation bundle \(I\),
specifications \(S=\{s_1,\ldots,s_n\}\), imports \(M\), and an editable
slot schedule \(E\). In proof mode, the target is the fixed canonical
implementation \(I_c\) and each accepted slot establishes either
\(s_i(I_c)\) or its permitted audit polarity. In codeproof mode, the
candidate chooses the editable implementation \(I\) and must prove the
specifications or submit one of Vero's exact audit forms.

A repository passes only when every required specification passes the
independent grader. This all-or-nothing repository metric means local
proof coverage and global completion are materially different
objectives. A system that independently solves easy \(s_i\) may score
well by specification while leaving most repositories incomplete.

\subsection{2.2 Authority}\label{authority}

Our governing rule is:

\begin{Shaded}
\begin{Highlighting}[]
\NormalTok{model = untrusted proposer}
\NormalTok{deterministic/formal validator = authority}
\end{Highlighting}
\end{Shaded}

Every candidate follows:

\begin{Shaded}
\begin{Highlighting}[]
\NormalTok{pinned {-}\textgreater{} built {-}\textgreater{} checked {-}\textgreater{} admitted or refuted}
\end{Highlighting}
\end{Shaded}

The model cannot declare a theorem proven, change the allowed edit
surface, waive a residual, introduce a new axiom, or define the
benchmark result. Similarly, a Synthesa receipt proves identity,
authorization, and lineage; it does not independently reproduce Lean
semantics. Vero's own clean render, build, axiom screening, and report
remain the external score authority.

\subsection{2.3 Boundedness}\label{boundedness}

Benchmarks must expose a finite set of repositories, declarations,
slots, and grading outcomes. That does not make every theorem decidable,
but it creates many bounded control problems around theorem
construction. Synthesa compiles those control decisions: obligation
identity, dependency coverage, candidate selection, tactic budgets,
adapter completeness, certificate coverage, slot authorization, and
result admission. Mathematical cores are handled by Lean,
SMT/SyGuS-backed synthesis where appropriate, finite reflected
certificates, or a narrowly scoped proposal. The product never equates a
bounded workflow with a proof of an unbounded mathematical proposition.

\section{3. Compiled proof-guidance
architecture}\label{compiled-proof-guidance-architecture}

\subsection{3.1 Elaboration-first obligation
extraction}\label{elaboration-first-obligation-extraction}

Early experiments showed that source syntax is an unsafe proof-planning
interface. Names can be qualified differently after elaboration;
implicit arguments, coercions, typeclass instances, and reducible
definitions can change the actual goal. The system therefore queries
Lean's environment for every \texttt{spec\_*} declaration and
serializes:

\begin{itemize}
\tightlist
\item
  fully qualified declaration name and universe parameters;
\item
  all explicit and implicit binders with elaborated types;
\item
  the target expression;
\item
  referenced declarations and implementation heads;
\item
  imports, source location, and dependency neighborhood;
\item
  render, toolchain, and elaboration-environment digests.
\end{itemize}

An obligation is eligible for proof construction only if the declaration
resolves, its type elaborates, its implementation binding is identified,
and the record remains stable across clean renders. No source-text
fallback is allowed. This converts a generic build failure into a
precise distinction between a malformed input, an unknown declaration,
and a genuine proof goal.

\subsection{3.2 Parent contracts and GIGO
protection}\label{parent-contracts-and-gigo-protection}

A reusable theorem needs an explicit parent world: types, operations,
laws, and representation bridges sufficient for its statement and proof.
We compile finite adequacy surfaces whose cells answer whether each
required clause is bound in the frozen source, needs a bridge lemma, is
inapplicable, or is missing. Unknown cells default to
\texttt{GIGO\_HOLD}.

This gate prevents a dangerous failure mode. When an intended law is
absent, repeated tactics or a powerful model may produce increasingly
plausible text, but no axiom-clean proof follows. The system instead
emits a minimal residual such as
\texttt{MISSING\_EDGE\_ORIENTATION\_BRIDGE},
\texttt{MISSING\_PROGRESS\_MEASURE}, or
\texttt{BENCHMARK\_SPEC\_DOES\_NOT\_ELABORATE}. A missing assumption is
not silently inserted into the theorem.

GIGO is distinct from Vero's formal audit. A formal audit requires an
accepted Lean proof that a fixed reference violates a spec, a spec is
unsatisfiable, or a set is inconsistent. Under-specification or opacity
may establish neither polarity and therefore remains unfilled.

\subsection{3.3 Deterministic proof floor and residual
normalization}\label{deterministic-proof-floor-and-residual-normalization}

Before any model call, a frozen portfolio of Lean tactics and typed
proof constructors is tried under bounded time and resource limits. The
portfolio is intentionally stable across a qualification stage so a
change in coverage can be attributed to representation or knowledge
rather than tactic shopping.

Rejected candidates are normalized into structured residuals, including:

\begin{Shaded}
\begin{Highlighting}[]
\NormalTok{UNRESOLVED\_GOAL          MISSING\_REWRITE}
\NormalTok{MISSING\_CASE\_SPLIT       INDUCTION\_REQUIRED}
\NormalTok{TYPE\_MISMATCH            MISSING\_INSTANCE}
\NormalTok{ARITHMETIC\_RESIDUAL      UNKNOWN\_DECLARATION}
\NormalTok{TERMINATION\_ISSUE        DEPENDENCY\_CYCLE}
\end{Highlighting}
\end{Shaded}

Normalization retains the exact Lean diagnostic and a digest. The class
is a routing signal, not a replacement for the diagnostic. It supports
deterministic selection of the next constructor, library query, adequacy
check, or proposal type.

\subsection{3.4 Semantic obligation
graph}\label{semantic-obligation-graph}

The repository is compiled into a graph with nodes for specifications,
implementation APIs, elaborated declarations, residuals, candidate
abstractions, verified lemmas, and module/import units. Edges record
reference, implementation binding, import, proof dependency, candidate
coverage, and verified discharge.

Clustering operates on semantic features rather than filenames alone:

\begin{itemize}
\tightlist
\item
  target head and normalized proposition shape;
\item
  binder type heads and shared state types;
\item
  referenced definitions and recursive predicates;
\item
  implementation head;
\item
  dependency neighborhood; and
\item
  failed tactic/residual signature.
\end{itemize}

The output is a bounded assertion such as ``these nine unresolved
obligations share state \(X\), transition \(T\), and preservation-shaped
targets,'' not a claim that any particular invariant is true. That
distinction leaves abstraction invention open while compiling the
repository-organization problem.

\subsection{3.5 Oracle and certificate
compilation}\label{oracle-and-certificate-compilation}

We use ``oracle'' to mean a pinned, checkable knowledge artifact, not an
authority that can override Lean. Candidate sources include:

\begin{itemize}
\tightlist
\item
  existing Lean libraries;
\item
  public formal source documents pinned by URL, revision, license, and
  digest;
\item
  finite truth tables or certificates for bounded domains;
\item
  CEGIS/SMT/SyGuS-produced candidates with checkable witnesses;
\item
  previously verified repository helpers; and
\item
  one narrowly typed model proposal.
\end{itemize}

Every source-derived theorem needs a repository adapter. A public Coq
theorem, for example, can explain intended FLoCq semantics but does not
prove a proposition about a differently translated opaque Lean constant.
The adapter must establish the representation and operation equations in
the frozen environment. Adapter adequacy is a separate, fail-closed
seam.

Finite certificates can be atomized. Instead of trusting a generated
table as a monolith, the system binds the domain, one certificate row
per input, the decision rule, coverage proof, and projection theorem.
This was useful for finite-field reasoning: computation supplies bounded
evidence while Lean checks how it implies each specification.

\subsection{3.6 Verified lemma library}\label{verified-lemma-library}

An accepted helper is stored with exact source bytes, proposition,
imports, dependency roots, Lean/toolchain/environment digest, axiom
report, covered specifications, goal-shape tags, and verification
receipt. The library admits only Lean-checked artifacts. Dependencies
form an explicit acyclic graph and are topologically materialized.

The important unit of progress is therefore not a transient model
response but a verified abstraction that can discharge several
specifications---or several repositories---without further generation.
The system rechecks the library before proposing new work.

\subsection{3.7 Bounded proposer ABI}\label{bounded-proposer-abi}

The proposer is invoked only after deterministic paths end at one exact
semantic residual. A request contains:

\begin{itemize}
\tightlist
\item
  proposal class, such as helper lemma, invariant, proof skeleton, or
  proof-friendly implementation;
\item
  exact elaborated binders and target;
\item
  permitted imports and output grammar;
\item
  verified local context and counterexamples;
\item
  accepted and rejected examples;
\item
  the smallest residual to repair; and
\item
  repository, Vero, Lean, prompt, model, and configuration digests.
\end{itemize}

The response bytes are pinned before validation. A rejected response may
produce a narrower residual but never mutates the verified library. If a
proposal is admitted, the resulting theorem---not the model
response---is retained as deterministic knowledge. Identical future
residuals should reuse it with zero new calls.

\subsection{3.8 Signed Act and independent
grading}\label{signed-act-and-independent-grading}

Candidate materialization remains in Synthesa's existing actuator path:

\begin{Shaded}
\begin{Highlighting}[]
\NormalTok{software.plan {-}\textgreater{} signed plan {-}\textgreater{} software.build {-}\textgreater{} terminal receipt}
\end{Highlighting}
\end{Shaded}

The plan binds the benchmark, mode, source and sandbox digests, Vero and
grader digests, Lean/Lake/Python environment, exact allowed slots,
parent-world root, abstraction/library roots, builder binary, and
terminal-once state. Build rechecks these values and refuses drift,
tamper, unauthorized slots, and replay.

Finally, Vero extracts only scheduled slot bodies, re-renders from the
frozen benchmark, drops unexpected markers, builds the clean project,
and performs its axiom checks. Only that report is aggregated.

\section{4. Experimental method}\label{experimental-method}

\subsection{4.1 Environment}\label{environment}

We froze Vero at commit
\texttt{0a7325df9e9e6dbc275c0ad483b3d1cbe38d9b09} and used
\texttt{leanprover/lean4:v4.29.1}. The final campaign compiler has
SHA-256
\texttt{0653c7421f2aa681fb8d7795b400fc43cec2e3cfd24be85ecc928e921f2666aa}.
The aggregate campaign root is
\texttt{cfd0bc418a41ca32eb5b8ede967d2e55c4be38ca514fc0b6e503369711ea3d8c}.

The study was iterative and benchmark-seen. Capabilities were developed
in response to observed residual classes, then frozen and checked
against further repositories where possible. This is product-oriented
research, not a blind held-out agent evaluation. We preserve strict
failures and distinguish generic compiler additions from repository
proof artifacts, but the absence of a prospective holdout limits causal
conclusions.

\subsection{4.2 Starting and stopping
points}\label{starting-and-stopping-points}

The campaign starts from a frozen 43-repository portfolio with 2
complete repositories and 1030 accepted specifications. Subsequent
results are promoted only when backed by an official report and, for the
qualified Act path, a signed plan and receipt. We stopped at 40/43 and
2606/2705 rather than inventing parent laws for the final 99
obligations.

\subsection{4.3 Modes}\label{modes}

The selected final artifact is proof mode for 37 repositories and
codeproof for six. This allows the systems study to use the strongest
qualified artifact per repository, but it prevents direct comparison
with one Vero mode. We therefore report three views: the mixed
portfolio, the 37-repository proof subset, and the six-repository
codeproof subset. A future leaderboard experiment must use a
maintainer-approved uniform protocol.

\subsection{4.4 Metrics}\label{metrics}

Primary measurements are complete repositories, accepted specifications,
and false acceptance. Secondary measurements include proposal calls and
cost, verified abstractions, residual classes, specs discharged per
helper, and repository unlocks per generic abstraction. The last metric
tests the central industrialization hypothesis: if each new abstraction
unlocks a repository, the method scales; if every tail repository needs
bespoke theory, it does not.

\subsection{4.5 Cost accounting}\label{cost-accounting}

The evidence tree contains 33 uniquely priced NanoGPT proposer calls:
628,022 prompt tokens, 34,377 completion tokens, 662,399 total, and USD
0.710391762755. The models were Qwen 3.5 4B, DeepSeek V4 Flash, GPT-5.5,
and Claude Sonnet 4.6. This conservative ledger includes rejected and
diagnostic calls.

Frozen artifact grading uses no model. ``Zero model calls'' means zero
calls in replay because accepted knowledge has been compiled into the
artifact; it does not erase the 33 development calls. The dollar ledger
also excludes Codex orchestration, human labor, and local compute and
therefore is not total system cost.

\section{5. Results}\label{results}

\subsection{5.1 Aggregate outcome}\label{aggregate-outcome}

\begin{longtable}[]{@{}lrrrr@{}}
\toprule\noalign{}
Result view & Repositories & Full & Specifications & Spec rate \\
\midrule\noalign{}
\endhead
\bottomrule\noalign{}
\endlastfoot
Starting portfolio & 43 & 2 & 1030/2705 & 38.08\% \\
Final mixed-mode portfolio & 43 & 40 & 2606/2705 & 96.34\% \\
Selected proof subset & 37 & 35 & 2286/2382 & 95.97\% \\
Selected codeproof subset & 6 & 5 & 320/323 & 99.07\% \\
\end{longtable}

The portfolio adds 38 complete repositories and 1,576 accepted
specifications relative to its starting state. All 43 final artifacts
were extracted using upstream \texttt{vero-extract}; corresponding
report and artifact digests are provided for independent maintainer
grading.

The aggregate is strong evidence that deterministic organization and
reusable verified knowledge can materially improve a repository-scale
proof campaign. It is not evidence that the same frozen system would
score 40/43 on a new blind sample, nor that it beats 27/43 under Vero's
published single-mode budget.

\subsection{5.2 Ntheory: one abstraction chain, five residual
projections}\label{ntheory-one-abstraction-chain-five-residual-projections}

Ntheory moved from 57/62 to 62/62 through a verified number-theory
chain: prime factorization support, a character contract,
Gauss/Eisenstein and quadratic reciprocity support, an abstract Jacobi
factor product, and a refinement from the frozen recursive worker to
that abstraction. The final foundation was projected into five residual
specifications. The closing execution made no model call, though earlier
R\&D proposer calls are included in the global ledger. This case
demonstrates amortization within a hard repository cluster.

\subsection{5.3 Reedsolo: bounded computation plus algebraic
seams}\label{reedsolo-bounded-computation-plus-algebraic-seams}

Reedsolo reached 48/48 in proof mode through a staged finite-field and
polynomial-semantics program. A complete byte-domain certificate
supplied finite evidence; Lean-checked field laws, polynomial
evaluation, synthetic division, generator recurrence, root, and distance
abstractions connected it to repository specifications. The key lesson
is that finite enumeration alone was insufficient: its certificate
needed atomized coverage and verified seams to the higher-level
propositions.

\subsection{5.4 Shared helpers as repository
infrastructure}\label{shared-helpers-as-repository-infrastructure}

Across the campaign, difficult walls were approached by compiling shared
parents before exact \texttt{prove\_*} projections. Examples include
ordered-map set-update laws, relation/path composition, codec
roundtrips, periodic calendar oracles, order/canonicalization
properties, and state-transition invariants. Successful parents became
dependency-addressed inputs rather than prompt text. This qualitatively
aligns with Vero's observation that completed repositories are organized
around shared helper libraries (\citeproc{ref-vero2026}{Ye, Lou, et al.
2026}).

\subsection{5.5 Honest residuals}\label{honest-residuals}

The final 99 specifications are concentrated in three repositories:

\begin{longtable}[]{@{}
  >{\raggedright\arraybackslash}p{(\columnwidth - 8\tabcolsep) * \real{0.1765}}
  >{\raggedright\arraybackslash}p{(\columnwidth - 8\tabcolsep) * \real{0.1765}}
  >{\raggedleft\arraybackslash}p{(\columnwidth - 8\tabcolsep) * \real{0.2353}}
  >{\raggedleft\arraybackslash}p{(\columnwidth - 8\tabcolsep) * \real{0.2353}}
  >{\raggedright\arraybackslash}p{(\columnwidth - 8\tabcolsep) * \real{0.1765}}@{}}
\toprule\noalign{}
\begin{minipage}[b]{\linewidth}\raggedright
Repository
\end{minipage} & \begin{minipage}[b]{\linewidth}\raggedright
Mode
\end{minipage} & \begin{minipage}[b]{\linewidth}\raggedleft
Accepted
\end{minipage} & \begin{minipage}[b]{\linewidth}\raggedleft
Residual
\end{minipage} & \begin{minipage}[b]{\linewidth}\raggedright
Main finding
\end{minipage} \\
\midrule\noalign{}
\endhead
\bottomrule\noalign{}
\endlastfoot
DedekindReals & proof & 0/82 & 82 & frozen spec module has a dependent
type mismatch \\
FLoCq & proof & 189/203 & 14 & disconnected trusted wrappers and opaque
inverse contracts \\
Verdict & codeproof & 116/119 & 3 & opaque non-editable conversion lacks
preservation equations \\
\end{longtable}

These are not all equivalent. DedekindReals fails before obligations can
be faithfully elaborated. FLoCq includes two other specifications
already accepted through formal disproof, one translation-premise
concern, twelve wrapper alignment gaps, and one native-Float inverse
gap. Verdict's exposed conversion type admits both satisfying and
violating models, but neither polarity has been proved for the frozen
opaque constant. Stopping preserves the distinction between ``cannot
prove,'' ``input is malformed,'' and ``formal audit accepted.''

\section{6. Analysis}\label{analysis}

\subsection{6.1 Representation precedes proof
search}\label{representation-precedes-proof-search}

The first deterministic proof-floor attempt originally produced zero
coverage because it did not operate on the true elaborated
\texttt{spec\_*} types. Expanding the tactic portfolio would have
confounded the experiment. Resolving declarations, binders, targets, and
implementation heads first converted build noise into meaningful proof
residuals. This is a general lesson for proof agents: compiler feedback
is only useful when attributed to a faithful semantic unit.

\subsection{6.2 Compile organization, generate only
abstractions}\label{compile-organization-generate-only-abstractions}

The model is poorly positioned to repeatedly rediscover repository
structure. It is better used for the irreducible question ``what lemma
or invariant might connect these exact obligations?'' after the system
has compiled the cluster, known context, counterexamples, and output
shape. This reduces authority and scope at once. The accepted result can
then be reused without repeating the call.

\subsection{6.3 Precompiled public documents are useful but not
magical}\label{precompiled-public-documents-are-useful-but-not-magical}

Public formalizations, manuals, standards, and reference implementations
can be compiled into candidate oracles. They accelerate theorem
discovery and clarify intended semantics. They do not automatically
cross a representation seam. FLoCq and Verdict show why: a public
theorem about a valid Coq type or a Rust conversion cannot prove a
theorem about an opaque or differently modeled Lean declaration until a
checked adapter exists.

\subsection{6.4 GIGO is a product
capability}\label{gigo-is-a-product-capability}

Failing closed on an under-specified parent may lower a nominal score,
but it is essential for trustworthy automation. The same decision
machinery used for software policies can compile the completeness
surface of a proof adapter. Every declared requirement must have an
explicit disposition; missing cells produce exact residuals. This
prevents benchmark pressure from becoming a reason to manufacture
assumptions.

\subsection{6.5 What zero-call replay
buys}\label{what-zero-call-replay-buys}

Zero-call replay changes the economics and governance of verified
artifacts. Once a lemma or finite certificate is admitted, future builds
can use its content-addressed proof without provider availability,
stochastic variation, or inference cost. An organization can audit which
knowledge was introduced, which obligations it covers, and which
compiler/kernel accepted it. This is amortization, not retroactive model
erasure.

\section{7. Relationship to prior
work}\label{relationship-to-prior-work}

Vero establishes repository-scale joint implementation and proof
synthesis as a benchmark and shows that shared helper organization
separates many full solves from partial ones
(\citeproc{ref-vero2026}{Ye, Lou, et al. 2026}). We use Vero's existing
authority boundary and audit mechanism rather than replacing them, and
focus on compiling the organization layer its evaluation identifies.

Verina studies joint generation of code, specifications, and proofs at a
smaller task granularity and shows that iterative Lean feedback improves
proof success but remains costly and limited even over many repair steps
(\citeproc{ref-verina2025}{Ye et al. 2025}). Our residual normalizer has
a related motivation, but narrows repair to elaborated repository
obligations and retains accepted results as verified shared
infrastructure.

VeriSpecGen decomposes natural-language intent into atomic requirements,
associates tests with explicit traceability, and performs localized
repair (\citeproc{ref-verispecgen2026}{Ye, Yang, et al. 2026}). Synthesa
applies an analogous decomposition principle on the other side of the
specification boundary: formal declarations are elaborated, clustered,
and linked to proof residuals and reusable abstractions. We do not
synthesize Vero's frozen specifications, and Lean proof validity is
stronger than test-based intent evidence. The works are complementary.

CEGIS and SyGuS treat synthesis as bounded candidate generation refined
by counterexamples and formal constraints
(\citeproc{ref-cegis2006}{Solar-Lezama et al. 2006};
\citeproc{ref-sygus2013}{Alur et al. 2013}). Synthesa uses this style
where domains and grammars are suitable, while allowing a model to
propose abstractions outside the deterministic grammar. In both cases
the candidate is subordinate to the verifier.

\section{8. Business and product
implications}\label{business-and-product-implications}

The business value is not merely a higher benchmark number. The
architecture turns one-time expert or model insight into governed
reusable capital. A verified helper or compiled standards oracle can be
reused across builds with no inference call, while lineage answers
exactly which artifact, environment, and authority produced a result.
This is attractive for regulated and high-assurance software where
repeatability, reviewability, and refusal are as important as generation
speed.

The workflow also separates expensive uncertainty from routine
execution. Deterministic analysis handles obligation discovery,
coverage, policy, dependency scheduling, and admission. Models are
reserved for missing abstractions. Lean and domain validators settle
truth. Signed Act constrains mutation. Independent graders or enterprise
postconditions settle deployment success. This pattern extends beyond
Lean to quality, security, compliance, and transformation domains when
their decisions can be compiled into bounded oracle contracts with
checkable seams.

Finally, the residual repositories illustrate commercial value in saying
no. An automated system that invents a contract for an opaque conversion
can appear productive while creating unreviewable risk. A precise
hold---``the required preservation equation is absent from the frozen
parent''---is an actionable engineering output. It directs a spec owner
to the smallest repair and avoids spending more model budget on an
impossible request.

\section{9. Limitations and threats to
validity}\label{limitations-and-threats-to-validity}

First, this was an iterative engineering campaign with access to
benchmark repositories and residual reports. Results may include
benchmark-specific proof artifacts even though compiler logic is
intended to remain generic. A blind holdout is needed to measure
generalization.

Second, the final portfolio mixes modes. The proof and codeproof subsets
cover different repositories, so their completion rates cannot be
extrapolated to all 43. A leaderboard comparison requires a uniform mode
or explicit organizer approval of portfolio selection.

Third, cost accounting is incomplete. The NanoGPT ledger is exact for
the 33 captured calls, but orchestration, labor, and compute are not
priced. Frozen replay cost is operationally useful but is not comparable
to one-pass agent development cost.

Fourth, zero false acceptance is measured relative to the campaign's
official-report admission rule. It is not a proof that every
specification captures intended source semantics. Indeed, the curator
findings expose translation and parent-contract risks that only
benchmark maintainers can resolve.

Fifth, source availability is split. The extracted artifacts and reports
are prepared for private maintainer replication. The research compiler
currently lives on a preserved dirty worktree rather than a clean
release commit. A public research artifact should freeze and license a
reproducible source snapshot or provide a minimal open agent wrapper.

Sixth, the campaign did not quantify total helper-line share,
cross-repository reuse, or repository-unlocks-per-abstraction uniformly
across all stages. Future runs should emit those metrics directly from
signed artifact graphs.

\section{10. Next experiment}\label{next-experiment}

The next scientifically clean experiment is not more work on the same 99
residual specifications. It is a frozen, uniform-mode, prospective
campaign:

\begin{enumerate}
\def\labelenumi{\arabic{enumi}.}
\tightlist
\item
  agree with Vero maintainers on proof or codeproof mode, time, compute,
  and cost accounting;
\item
  freeze the compiler source, binary, tactic portfolio, oracle library,
  proposer ABI, prompts, model versions, and public-document corpus;
\item
  select an unseen repository subset or newly curated Vero release;
\item
  run end to end with complete event, cost, residual, and lineage
  capture;
\item
  submit only \texttt{vero-extract} artifacts for independent grading;
  and
\item
  report ablations for deterministic floor, compiled-oracle reuse, small
  proposer, and stronger proposer.
\end{enumerate}

The key causal metric should be new full repositories per newly admitted
generic abstraction, alongside model calls per accepted specification. A
positive cross-repository ratio with zero false acceptance would
directly test whether compiled organization generalizes.

\section{11. Conclusion}\label{conclusion}

Vero asks whether agents can build formally verified repositories. Our
study suggests a complementary question: how little authority and
repository organization must be left to the agent? Synthesa compiles
elaborated obligations, deterministic search, semantic clusters, parent
adequacy, certificates, verified helper dependencies, and signed
execution into a fail-closed workflow. Models propose only bounded
missing abstractions; Lean decides proof validity; Vero decides the
score.

The resulting 40/43, 2606/2705 mixed-mode portfolio is encouraging but
must be read with its experimental qualifications. More important than
the headline is the observed mechanism: verified abstractions can break
multi-specification walls and become zero-call reusable infrastructure,
while malformed or under-specified parents stop with exact residuals.
That is a plausible product path from stochastic code generation toward
governed verified synthesis.

\section{Data and artifact
availability}\label{data-and-artifact-availability}

The full per-repository results, cost ledger, issue drafts, artifact
hashes, and a private 43-artifact \texttt{vero-extract} bundle accompany
this draft. The extracted artifacts will be shared privately with the
Vero maintainers for independent grading and will not be made public
unless requested. Author list, company affiliation, artifact license,
source-release boundary, and AI-use disclosure must be confirmed before
archival submission.

\phantomsection\label{refs}
\begin{CSLReferences}{1}{0}
\bibitem[\citeproctext]{ref-sygus2013}
Alur, Rajeev, Rastislav Bodik, Garvit Juniwal, Milo M. K. Martin, Mukund
Raghothaman, Sanjit A. Seshia, Rishabh Singh, Armando Solar-Lezama,
Emina Torlak, and Abhishek Udupa. 2013. {``Syntax-Guided Synthesis.''}
In \emph{Proceedings of FMCAD}.

\bibitem[\citeproctext]{ref-cegis2006}
Solar-Lezama, Armando, Liviu Tancau, Rastislav Bodik, Sanjit Seshia, and
Vijay Saraswat. 2006. {``Combinatorial Sketching for Finite Programs.''}
In \emph{Proceedings of ASPLOS}.

\bibitem[\citeproctext]{ref-vero2026}
Ye, Zhe, Hantao Lou, Yuechun Sun, Peiyang Song, Zhengxu Yan, Timothe
Kasriel, Qingyang Zhang, et al. 2026. {``Vero: Can AI Agents Build
Formally Verified Software Repositories?''} \emph{arXiv Preprint
arXiv:2608.13522}. \url{https://arxiv.org/abs/2608.13522}.

\bibitem[\citeproctext]{ref-verina2025}
Ye, Zhe, Zhengxu Yan, Jingxuan He, Timothe Kasriel, Kaiyu Yang, and Dawn
Song. 2025. {``Verina: Benchmarking Verifiable Code Generation.''}
\emph{arXiv Preprint arXiv:2505.23135}.
\url{https://arxiv.org/abs/2505.23135}.

\bibitem[\citeproctext]{ref-verispecgen2026}
Ye, Zhe, Aidan Z. H. Yang, Huangyuan Su, Zhenyu Liao, Samuel Tenka,
Zhizhen Qin, Udaya Ghai, Dawn Song, and Soonho Kong. 2026.
{``Intent-Aligned Formal Specification Synthesis via Traceable
Refinement.''} \emph{arXiv Preprint arXiv:2604.10392}.
\url{https://arxiv.org/abs/2604.10392}.

\end{CSLReferences}

\end{document}
